\section{Conclusion}
\label{sec:conclusion}

This paper studies $L_0$ regularization as a way to subsample a large microsimulation dataset down to
a usable size. The problem arises because a pipeline built for fidelity combines many sources and
adds record-level variation, which produces more candidate records than a deployable dataset can
hold. The reduction to a fixed budget is a sampling choice, and the paper treats it as one.

The method uses Hard Concrete gates to select records and fit their weights in one optimization, with
selection driven by the calibration targets. That makes the retained-record count an explicit control
and keeps the reduced dataset calibrated to its targets. We compare the method against two
matched-budget baselines that share its calibrator, random sampling with reweighting and survey-weight
(Hansen--Hurwitz) sampling, and report how the size, geography, and weight controls behave.

The contribution is therefore twofold: a target-informed sampling method, and an open pipeline,
\populace{}, that makes the method a configuration choice other projects can reuse. Future work
can benchmark the sampler against classical calibration estimators and convex sparse methods,
formalize the out-of-sample target evaluation, and extend the geographic detail of the candidate
universe.
